\chapter{Conclusión}
El presente Informe Final de Seminario de Título consolida la fundamentación metodológica, el estado del arte, el diseño y una primera implementación de un sistema probabilístico basado en Cadenas Ocultas de Markov (HMM) para la detección temprana de Amenazas Persistentes Avanzadas (APT).

La Revisión Sistemática de la Literatura, conducida bajo las directrices de Kitchenham y documentada según la declaración PRISMA 2020, analizó un corpus de 77 estudios primarios de Scopus, Web of Science e IEEE Xplore y evidenció cuatro brechas: la dependencia de alertas previamente correlacionadas, la escasa validación de modelos HMM sobre flujos de tráfico, la limitada representatividad temporal de los conjuntos de datos y el predominio de métricas globales como \textit{Accuracy}, que enmascaran el desempeño frente a las clases minoritarias.

Para responder a estos vacíos, la propuesta adoptó la metodología Design Science Research (DSR) y formuló un HMM generativo que vincula directamente el tráfico de red con las etapas del ataque, evaluado sobre DAPT2020 y Unraveled para asegurar tanto la comparabilidad con el estado del arte como la validez ecológica sobre trazas de varias semanas.

El prototipo implementado (Capítulo~\ref{cap:experimentos}) confirma que la estimación supervisada mantiene la correspondencia entre estados y etapas que el entrenamiento no supervisado pierde, y que la verificación posterior impide que el refinamiento con Baum-Welch la degrade. A su vez, muestra que con la configuración actual los modelos de referencia superan al HMM en la clasificación por ventana y en la detección a igual tasa de falsas alarmas.

El trabajo futuro se concentrará en:
\begin{enumerate}
    \item Incorporar una línea de tiempo de la campaña a nivel de red y un atributo de volumen saliente acumulado para la exfiltración.
    \item Seleccionar y congelar en validación la configuración del preprocesamiento, incluida la combinación de la ventana de un minuto con la discretización por terciles.
    \item Contrastar las estrategias frente al desbalance que preservan la estructura temporal de las secuencias, reservando SMOTE para los modelos de referencia.
    \item Evaluar el modelo por etapa y en detección temprana frente a la tasa de falsas alarmas, con validación cruzada por día, y compararlo con Random Forest y XGBoost bajo las mismas condiciones.
\end{enumerate}

Con ello, se busca proporcionar un artefacto funcional que reduzca el tiempo de exposición frente a campañas dirigidas en infraestructuras críticas.
